\scanpage{47}
\uline{\unsure{Consequence|Consequences}}: If $\bigcup_{A\in\Sigma}A^*=S$, then \deleted{the} $(A^*)_{A\in\Sigma}$ is an \deleted{\unsure{of the|without|in}} \unsure{un}mixed partition of $S$. If $\bigcup_{A\in\Sigma}A^*\ne S$, then setting
\[
 \deleted{\unsure{R,\Sigma|R\ \Sigma}}\mathpunct{.}\qquad
 \left\{\begin{aligned}
 R&=S\setminus\bigcup_{A\in\Sigma}A^*,\\
 R^*&=R\setminus\bigcup_{\substack{A\in\widehat\Sigma\\ A\subsetneq R}}A=R
 \end{aligned}\right.\qquad
 \begin{aligned}
 \widehat\Sigma&=\Sigma\cup\{R\}\\
 &\deleted{\unsure{\Sigma|\in\Sigma|\ne\Sigma_0}}
 \end{aligned}
\]
then $(A^*)_{A\in\widehat\Sigma}$ is an unmixed partition of $S$.

\uline{NB} If $R$ is itself a segment, then \deleted{\unsure{the interest} \unsure{of taking it|of adjoining it}} $\widehat\Sigma$ satisfies the conditions of $\Sigma$ -- and $R$ is an \deleted{``\unsure{sp.|spéc.}} ``isolated'' element of $\widehat\Sigma$, in the sense that \inserted{for} \unsure{every|$A$|$\forall$} other \unsure{element|element} \inserted{$A$} of $\widehat\Sigma$, \deleted{\unsure{one}}. $R\cap A=\phi$ -- \deleted{\unsure{therefore}} \inserted{i.e.} $R$ is \unsure{:|at} once minimal, and maximal. Consider
\[
 \begin{aligned}
 (\sigma\ 3)\qquad R=S\setminus\bigcup_{A\in\Sigma}A\;
 \bigl(=S\setminus\bigcup A^*\bigr)&\quad\text{is}\\
 &\quad\text{not a segment}
 \end{aligned}
\]
(NB \deleted{\unsure{for example|p.}}\inserted{\unsure{this is}} the case if $R=\phi$ for example)

% slanted
\margin{false}
Let us show in this case that \deleted{\unsure{one has|is}} $\Sigma$ is \deleted{defined} \inserted{\unsure{known|known}} when one knows the partition $\widehat\Sigma^*=\{A^*\mid A\in\widehat\Sigma\}$

\[
 \bigl(\text{where }\widehat\Sigma=\left\{\begin{aligned}
 &\Sigma\quad\text{if }R=\phi\\
 &\Sigma\cup\{R\}\quad\text{if }R\ne\phi
 \end{aligned}\right.\bigr)
\]
% archive sheet number: 46
